\section{Roman Discipline Concerning the ``Secret''}

Roman action at La Salette is often summarized in two incompatible slogans: ``Rome condemned La Salette'' or ``Rome never acted against the secret.'' Both are false because each suppresses the object. The Holy See did not revoke Grenoble's event judgment or forbid the authorized devotion. It did take concrete action against publication, discussion, and particular books associated with the expanded secret.

\subsection{The archival Holy Office measures of 1880}

The 1879 Lecce booklet quickly provoked controversy in France and Italy. Michel Corteville and René Laurentin later published and discussed an archival Holy Office dossier from 1880. On their account, an internal process led to practical measures against distribution of the booklet: copies already circulated were to be withdrawn as far as possible, and Mélanie was not to continue similar writing or explanation. Cardinal Prospero Caterini conveyed the action to relevant ecclesiastical superiors and bishops.

This was real reported Roman discipline, even though the internal material was not a public AAS text and the correspondence did not offer a proposition-by-proposition doctrinal analysis. This project did not inspect the archival originals, and the available secondary locus is too broad to make subsidiary dates, addressees, and decision mechanics carry independent argumentative weight. The intervention cannot fairly be described as Roman authentication of the booklet. It also should not be enlarged into a condemnation of Marian devotion, the mountain event, or the public message. The archival source and its scholarly publication status must travel with the claim.

\subsection{The universal decree of 21 December 1915}

The first controlling public Roman text in this sequence appears in \emph{Acta Apostolicae Sedis} 7 (1915), printed page 594, under the title \emph{Decretum circa vulgo dictum ``Secret de La Salette''}. The Holy Office notes that some persons, including clergy, continued to publish books, pamphlets, and periodical articles about the so-called secret, its different forms, and its application to present or future times, without or against their ordinaries' permission.

The decree commands all the faithful not to discuss or treat the subject under any pretext or form. It prescribes serious penalties then applicable to clergy and laity and invokes existing publication discipline. Its last sentence is indispensable:

\begin{quote}
\small\latin{Hoc autem decretum devotionem non vetat erga Beatissimam Virginem sub titulo Reconciliatricis vulgo de la Salette nuncupatam.}
\end{quote}

Thus the decree itself distinguishes secret controversy from devotion to the Blessed Virgin under the title of Reconciliatrix, commonly called La Salette. One cannot obey the first half by erasing the second, or honor the second by denying the first.

\subsection{The 1916 book action and its printed name variants}

On 12 April 1916 the Holy Office declared a two-volume work condemned and proscribed under the general rules. The AAS declaration at printed page 175 gives the title \emph{La Leçon de l'Hôpital Notre-Dame d'Ypres -- Exégèse du Secret de la Salette} and prints the author as ``Dr Henry Mariam.'' It identifies volume I as Paris, 1915, and volume II, \emph{Appendices}, as Montpellier, 1915.

The Congregation of the Index then listed the work in its decree of 5 June, printed at pages 178--179. There the author appears as ``Dr Henri Mariavé.'' The difference is in the official volume and must not be silently harmonized. The safest catalog statement is that the April declaration's ``Dr Henry Mariam'' and the June Index listing's ``Dr Henri Mariavé'' refer to the same titled work.

This action is narrower than the 1915 subject-wide command in one respect: it identifies a particular book. It does not name the 1846 public message as erroneous or undo authorized cult.

\subsection{The 1923 decree and the anomalous date}

On 9 May 1923 the Holy Office's general session proscribed and condemned a named booklet and directed that copies be withdrawn. Pius XI approved the resolution that day; the decree is dated 10 May. The operative bibliographic entry, printed in AAS 15 (1923), page 288, reads:

\begin{quote}
\small
\emph{L'apparition de la très Sainte Vierge sur la sainte montagne de la Salette le samedi 19 septembre 1845} [\emph{sic}]. --- \emph{Simple réimpression du texte intégral publié par Melanie, etc. Société Saint-Augustin, Paris-Rome-Bruges, 1922.}
\end{quote}

The historical event occurred in 1846; ``1845'' is the official entry's anomalous printed date and is preserved with \emph{sic}. Correcting it invisibly would corrupt the legal bibliography. The named object is the Société Saint-Augustin 1922 booklet describing itself as a simple reprint of Mélanie's full text. The decree therefore cannot reasonably be reduced to a generic warning against an unrelated appendix; neither should it be inflated into a condemnation of every sentence ever spoken at La Salette.

Later literature mentions additional Holy Office correspondence about the 1922 reprint. No sufficiently reproducible publication and page locus for that correspondence was established in this project, so it is not used to refine the decree's object. The controlling conclusion instead rests on the 1923 AAS entry itself, which names the Société Saint-Augustin booklet and describes it as a reprint of Mélanie's text.

\subsection{What changed in 1966 and 1983}

On 14 June 1966 the Congregation for the Doctrine of the Faith issued its notification on the status of the Index (AAS 58 [1966], 445). The Index no longer had the force of ecclesiastical law with attached censures, while the notification still appealed to the faithful's conscience under the demands of natural law. A decree of 15 November confirmed that the old automatic book prohibitions and penalties in canons 1399 and 2318 of the 1917 Code no longer bound as ecclesiastical law. The 1983 Code, canon 6 \S1, 3$^\circ$, then abrogated Apostolic See penal laws not retained in the new Code.

Those sources settle the old Index machinery and penalty schedule; they do not, without a further competent interpretation, settle the present juridic status of every standalone non-penal command in the 1915 decree. Consequently, this study neither presents the old censures as currently executable nor asserts that the subject-wide command remains binding today in its original breadth. It also does not call the 1879 or 1922 text ``rehabilitated.'' A change in legal machinery does not transform a proscribed book into an authenticated revelation, reverse the historical act, or require the Church to promote it. Current readers remain bound by faith, morals, ecclesial communion, truthfulness, and current competent discipline.

No current DDF act located through 16 July 2026 republishes the old penalties, authenticates the expanded secret, or reclassifies the 1851 event judgment. Questions about a concrete contemporary edition or cleric's public promotion belong to current competent authority, not private canonical improvisation.

\subsection{The acts compared}

\begin{fourstable}{Date and source}{Exact object}{Action}{Express or necessary limit}
1880; archival Holy Office dossier as published in scholarship & Mélanie's recently published booklet and further similar explanations & Reported withdrawal from circulation and restriction of similar writing or explanation, conveyed through Caterini & Archival originals not inspected; not a public AAS decree; no prejudice to the authorized La Salette cult.\\
21 Dec. 1915; AAS 7, 594 & Discussion of the so-called secret, its forms, and current/future applications & Universal prohibition with then-operative penalties & Expressly does not prohibit devotion to Mary as Reconciliatrix of La Salette.\\
12 Apr./5 June 1916; AAS 8, 175, 178--179 & \emph{La Leçon de l'Hôpital Notre-Dame d'Ypres} & Specific book condemned, proscribed, and listed & Author printed ``Henry Mariam'' and ``Henri Mariavé'' in the two notices; not an event decree.\\
9--10 May 1923; AAS 15, 287--288 & Named 1922 Société Saint-Augustin reprint of Mélanie's text & Booklet proscribed and condemned; copies withdrawn & Title prints 19 Sept. 1845 [\emph{sic}]; object is not the 1846 public message.\\
14 June 1966; AAS 58, 445 & Legal status of the Index generally & Ecclesiastical-law force and attached censures removed; conscience still addressed under natural law & Not a declaration that every formerly listed claim is true or safe, or a ruling on the current force of every separate earlier command.\\
\end{fourstable}

\interpretivekey{Object before slogan}{For every act ask: Who acted? On what date? On which text or practice? With what verb? Under what law? What was expressly preserved? At La Salette the answers yield a coherent pattern: positive reception of the 1846 event and cult, ecclesial appropriation of the public call to reconciliation, and sustained caution or restriction concerning the secret's expanded forms and apocalyptic application.}
